\chapter{Conclusiones}

El objetivo general de este proyecto de memoria fue desarrollar un servidor de nombres DNS iterativo en el lenguaje Rust, utilizando la librería \texttt{dns-rust}, y que cumpliera 
rigurosamente con el estándar de DNSSEC, siguiendo las especificaciones establecidas en los RFC 1034, 1035, 4033, 4034 y 4035. Se puede afirmar que dicho objetivo fue alcanzado de manera 
íntegra. El resultado final consistió en un proyecto de más de 3000 líneas de código puro, capaz de manejar consultas simultáneas tanto en UDP como en TCP, con soporte completo para EDNS0 y 
DNSSEC, y con la capacidad de gestionar múltiples zonas para emular adecuadamente la cadena de autenticación requerida por el modelo de seguridad DNSSEC.

Asimismo, se diseñó e implementó un algoritmo de búsqueda que responde de manera correcta y consistente para cada uno de los casos evaluados, siguiendo las reglas definidas por los estándares. 
El sistema incorpora firmas criptográficas reales que fueron generadas, insertadas y validadas correctamente durante el proceso de verificación. Además, cuenta con manejo explícito de errores, 
detección de loops y un desempeño estable incluso al procesar múltiples consultas en paralelo.

Durante el desarrollo también se obtuvieron conclusiones relevantes respecto a la estructura interna de los datos. En particular, se determinó que no es estrictamente necesario utilizar la 
estructura jerárquica de árbol sugerida por los RFC, y que incluso si hubiese sido deseable implementarla, Rust dificulta considerablemente este enfoque debido a su modelo de propiedad y los 
problemas que surgen al intentar crear estructuras de árboles mutables y compartidas. En consecuencia, se optó por un enfoque distinto, más adecuado para el lenguaje, que permitió avanzar sin 
comprometer el rendimiento ni la corrección del sistema.

Además, durante la implementación se identificaron varios problemas en la librería \texttt{dns-rust}, la cual se publicita como una librería lista y funcional. En concreto, se descubrieron 
dos bugs importantes: uno en las funciones de parsing y otro asociado al manejo de la sección \texttt{Additional}. También se evidenció la ausencia de funcionalidades fundamentales para un 
sistema DNSSEC, como la conversión de los distintos tipos \texttt{Rdata} a sus representaciones canónicas, lo que obligó a extender la librería para cumplir con los requisitos del protocolo.

A pesar de los avances, aún queda trabajo por realizar. No se logró implementar un sistema formal de testing utilizando herramientas especializadas, por lo que el proyecto actualmente carece 
de tests unitarios y de integración que permitan garantizar más exhaustivamente su correcto funcionamiento. Adicionalmente, el servidor está todavía lejos de ser una solución preparada para 
entornos reales, ya que faltan numerosos RFC por implementar, mecanismos de comunicación real con servidores externos, herramientas de seguridad como firewalls DNS para mitigar ataques, entre 
muchas otras funcionalidades necesarias para una operación a nivel productivo.

Sin embargo, a pesar de estas limitaciones, el proyecto constituye una base sólida y extensible para continuar el desarrollo de un servidor DNS iterativo con soporte DNSSEC completamente 
funcional. Se espera que esta implementación pueda servir como punto de partida para que el laboratorio \texttt{NIC Labs} continúe profundizando y expandiendo este trabajo en el futuro, 
si así lo estima conveniente.


